\section{The cubic theta normalization}
\label{app:whittaker}

This appendix supplies the Kazhdan--Patterson specialization and
normalization behind \Cref{prop:theta-input}.  We first specialize their
cover, inducing datum, and global representation, and derive the normalized
Whittaker factorization.  We then verify the exact lift identities and
compute the spherical Jacquet integral in the sections used in the paper.
The local calculation fixes the Gauss phase and the unit factor together,
and includes the modular factor in the triple-valuation step.  The main
argument uses the resulting global Fourier expansion with normalized local
factors and the exact local rules for the Weyl calculation.

\subsection{The selected cover and global representation}
\label{subsec:theta-specialization}

Denote the cubic Hilbert symbol at a place \(v\) by
\((\,\cdot\,,\,\cdot\,)_{3,v}\).  In the convention of Kazhdan and
Patterson, the restriction of their determinant-modified cocycle to the
diagonal torus is
\begin{equation}\label{eq:kp-torus-cocycle}
 \sigma_{c,v}\bigl(\operatorname{diag}(a,b),
                    \operatorname{diag}(a',b')\bigr)
 =(a,b')_{3,v}(ab,a'b')_{3,v}^{c},
 \qquad c\in\Z/3\Z.
\end{equation}
The root groups have their canonical splittings.  These conventions are
those of \cite[Section~0.1, pp.~39--44]{KP84}.

For the \(n\)-fold cover of \(\GL_d\), where \(d\) is the rank parameter,
the center criterion of \cite[Proposition~0.1.1, p.~42]{KP84} requires
\(z^{d-1+2dc}\in F_v^{\times n}\) for a scalar \(zI\).  For \(d=2,n=3\)
the exponent is \(1+4c\), which is zero modulo three exactly when \(c=2\).
Moreover, \eqref{eq:kp-torus-cocycle} on two scalars is
\[
 \sigma_{2,v}(zI,z'I)=(z,z')_{3,v}^{\,1+4\cdot2}=1.
\]
Thus their full preimage is central and the ordinary scalar section splits.
Quotienting its image gives the asserted local covers of \(\PGL_2\).
The adelic cover and its rational splitting are obtained by the reciprocity
construction in \cite[Section~0.2]{KP84}.  The ordinary scalar section
agrees with that splitting on rational scalars, so the quotient retains a
rational splitting.

We check the datum in \eqref{eq:theta-central-datum}.  A Hilbert symbol with
a cube argument is one; the cocycle on scalar pairs is also one.  Consequently
the cube and scalar sections commute and multiply without any remaining
cocycle.  If an element has two presentations in
\eqref{eq:theta-central-datum}, their two cube coordinates differ by a
common scalar cube, so the displayed absolute-value ratio agrees.  The
formula is therefore a genuine character, trivial on the split scalar
subgroup.  On a coroot cube it gives
\[
 \omega_v\bigl(s_v(\operatorname{diag}(t^3,t^{-3}))\bigr)=|t|_v.
\]
This is exactly the exceptional condition of
\cite[Section~I.2, p.~71]{KP84}.  The exceptional condition alone would
only specify the square of the value on \(a(\varpi^3)\); the datum in the
main text chooses the pure absolute-value value.

Use the idele norm in \eqref{eq:theta-central-datum} to define the global
character.  A rational diagonal lying in the adelic cube torus times scalars
has a rational ratio which is a cube at every completion.  If that ratio
were not a cube in \(F\), its nontrivial Kummer extension would have a
nonsplit unramified prime by Chebotarev.  Thus the ratio is a rational cube.
The product formula now shows that our character is trivial on the
intersection with the rational lift.  This is the compatibility hypothesis
in \cite[Section~II.1, p.~107, immediately after Lemma~II.1.1]{KP84}.
Extending by the trivial character on rational lifts gives the character
on the rational maximal abelian subgroup of that lemma.  As in
\cite[Section~II.1, pp.~107--109]{KP84}, one chooses an extension on the
product of local maximal abelian subgroups which agrees with the former
character on their intersection.  The character on the scalar center is
unitary.

At finite places, the local quotient and subrepresentation assertions are
\cite[Theorem~I.2.9(a),(b), p.~72]{KP84}.  At tame finite places, part (f)
of that theorem supplies the spherical vector for the normalized unramified
datum.  For a tame place, the uniqueness criterion is
\[
 d=n-1,\quad 2(c+1)=0\pmod n,
 \qquad\text{or}\qquad d=n
\]
by \cite[Corollary~I.3.6, p.~79]{KP84}.  It applies to
\((d,n,c)=(2,3,2)\).  The residue construction and restricted tensor
product are \cite[Theorem~II.2.1, p.~118]{KP84}.  Since the tame local
Whittaker spaces have dimension one,
\cite[Theorem~II.2.5(a), p.~122]{KP84} gives a nonzero global Whittaker
functional, and part (b) gives local uniqueness at every place.
At the complex places this also follows directly from
\cite[Theorem~I.6.5(c), p.~105]{KP84}.  The inducing spaces, their exceptional
images, and the automorphic residue all inherit the scalar action of
\(\omega\).  That action is trivial, so these representations descend to the
scalar quotients.

On the quotient torus the section of \(a(y)\) has cocycle
\((x,y)_{3,v}^{2}\).  Since \(-1\) is a cube in \(F\), this is an alternating
bimultiplicative law.  At a tame place it is trivial on pairs of units, and
its value on a uniformizer and a unit is the cubic residue symbol or its
inverse, with a common orientation at all places.  Choose the embedding
\(\iota\) with the orientation giving \eqref{eq:theta-effective-cocycle}.
The canonical compact splitting agrees with the sections of unit diagonals,
Weyl elements, and integral upper root elements by
\cite[Proposition~0.1.3, p.~44]{KP84}.  Conjugating by the Weyl element gives
the same agreement for the lower root group.
\cite[Corollary~I.3.4, p.~78]{KP84} gives relations for the coefficient
functions defining local Whittaker functionals.
\cite[Theorems~I.4.2--I.4.3, p.~87]{KP84} convert them to spherical
Whittaker values with the modular factor and give the central step.
The direct local calculation below proves
\eqref{eq:theta-local-values} and the lift identities
\eqref{eq:theta-lift-identities} in these conventions.

At a complex place the Hilbert symbol and the cover are trivial.
\cite[Theorem~I.6.4(c), p.~104]{KP84} identifies the exceptional
representation with the ordinary principal series when \(d\le n\).
Our cube datum fixes the characters on all of \(\C^\times\), because every
complex number has a cube root.  Triviality on scalars removes the determinant
character; equivalently any possible remaining finite-order continuous
character of the connected group \(\C^\times\) is trivial.  This gives the
characters in part \textup{(iii)} of \Cref{prop:theta-input}.  Their Bessel
function is computed in the main text from the Jacquet integral.  Finally,
the two Bruhat terms in the constant term of the Eisenstein series are
described by \cite[Proposition~II.1.2, p.~110]{KP84}.  At the exceptional
residue the surviving term belongs to the Weyl-conjugate induction, as in
\cite[Section~II.1, p.~114]{KP84}.  Its complex inducing character,
including the modular factor, is \(|y|_v^{1/3}\); in particular it has no
angular part.

\subsection{The normalized global Whittaker expansion}
\label{subsec:theta-global-normalization}

Let \(\lambda\) be the nonzero global functional just obtained.  The
restricted tensor product is spanned by pure tensors, so there is a pure
tensor \(\varphi^\circ\) with \(\lambda(\varphi^\circ)\ne0\).  Its
nondistinguished finite places form a finite set.  The final \(S\) contains
this set; the arithmetic data are constructed only after that final choice.
For \(v\notin S\), varying only the \(v\)-component of \(\varphi^\circ\)
gives a nonzero local Whittaker functional.  Normalize it at
\(\varphi_v^\circ\).  If other finitely many components have already been
varied, the resulting partial functional at \(v\) is still a Whittaker
functional, and local uniqueness makes it a scalar multiple of the
normalized one.  Evaluating at \(\varphi_v^\circ\) identifies that scalar.
Induction on the number of varied components proves
\eqref{eq:theta-normalized-factorization}.  The same argument factors the
complex components.

At a finite product of remaining places the functional descends to the
tensor product of the local twisted Jacquet quotients.  The twisted Jacquet
quotient of the principal series is finite dimensional by
\cite[Lemma~I.3.2, p.~75]{KP84}.  Since \(\Theta_v\) is its quotient by
Theorem~I.2.9(b), its twisted Jacquet quotient is a quotient of that
finite-dimensional space.  A functional on the finite tensor product is
therefore a finite sum of products.  This is the finite-S assertion in
\Cref{prop:theta-input}.

This normalization accounts for the correction to the original global
formula.  The erratum withdraws an auxiliary spherical-functional vanishing
assertion on pp.~119--120.  It replaces stabilization of the auxiliary
coefficient \(C_S\) by stabilization of \(C_ST(S)\), where in rank two
\[
 T(S)=\prod_{\substack{v\in S\\v\nmid\infty}}(1+Q_v^{-1}),
 \qquad Q_v=\#(\mathcal O_v/\mathfrak p_v).
\]
The sums in Theorems~II.2.2--II.2.3 become finite-\(S\) limits with the
factor \(T(S)^{-1}\), and the same change applies on p.~130
\cite[correction to Theorems~II.2.2--II.2.3]{KP85Erratum}.
The preceding argument begins with the nonzero automorphic functional and
normalizes each exterior local functional at the distinguished spherical
vector.  It proves the product \eqref{eq:theta-normalized-factorization}
without the withdrawn assertion or an unnormalized infinite product.

The rational splitting identifies every nonconstant Fourier coefficient
with this Whittaker function.  Conjugate \(n(x)\) by the rational lift of
\(a(\nu)\) and change \(x\) to \(\nu x\); the adelic Jacobian is one.
This gives \eqref{eq:theta-fourier-expansion}.  On a fixed compact set of
\(g\), smoothness at infinity and invariance under a fixed finite open
subgroup reduce the expansion along \(F\backslash\Adele_F\) to the smooth
Fourier expansion on a compact real torus.  Integration by parts in
sufficiently many archimedean directions gives absolute local uniform
convergence there.

\subsection{Sections and the spherical induced vector}

Fix an exterior finite place \(\mathfrak p\).  In this appendix write
\(\mathcal O=\mathcal O_{\mathfrak p}\), let \(\varpi\) be the chosen
uniformizer, put \(Q=q_{\mathfrak p}\), and abbreviate
\(\psi=\psi_{\mathfrak p}\), \(\chi=\chi_{\mathfrak p}\), and
\(\mathcal C=\mathcal C_{\mathfrak p}\).
Thus \(\psi\) has conductor \(\mathcal O\), additive measure gives
\(\mathcal O\) volume one, and \(\chi\) is the sextic residue
character on \(\mathcal O^\times\).  In all finite sums it is
extended by zero to the zero residue, including when its exponent is
zero.  The embedding of the cubic symbol has been chosen so that the
effective torus cocycle is
\begin{equation}\label{eq:app-torus-laws}
 \mathcal C(\varpi,u)=\chi(u)^{-2},\quad
 \mathcal C(u,\varpi)=\chi(u)^2,\quad
 \mathcal C(u,u')=1,\quad \mathcal C(\varpi,\varpi)=1
 \qquad(u,u'\in\mathcal O^\times).
\end{equation}
The last two identities follow from tameness and alternation.  In
particular pure uniformizer powers multiply without a cocycle.
All symbols with \(-1\) are one because \(-1\in F^{\times6}\).

We first record explicitly why the ordinary rank-one decompositions
are valid for these lifts.  See Kubota \cite[p.~114]{Kubota67} for the local
rank-one construction.  In the convention of
\cite[Section~0.1, p.~41]{KP84}, the determinant-modified formula is
\[
 \sigma_c(g_1,g_2)=
 \left(\frac{x(g_1g_2)}{x(g_1)},
       \frac{x(g_1g_2)}{x(g_2)\det g_1}\right)_3
       (\det g_1,\det g_2)_3^c,
 \quad
 x(g)=
 \begin{cases}
  g_{21},&g_{21}\ne0,\\
  g_{22},&g_{21}=0.
 \end{cases}
\]
We take \(c=2\) and then quotient the split scalar subgroup.
Substitution in this formula gives
\[
 n(x)a(y)=a(y)n(x/y),\qquad
 wa(y)=a(y^{-1})w
\]
on the quotient.  Before quotienting, the second equality has the
extra scalar \(yI\).  All its cocycles are symbols between powers of
\(y\) and \(-1\), hence are one.  For \(H\ne0\), the matrix identity
\[
 wn(H)=n(H^{-1})\operatorname{diag}(-H^{-1},H)\bar n(H^{-1})
\]
also lifts exactly.  The product of the first two factors has cocycle
one, and multiplication by the last factor has cocycle
\((H^{-1},-H)_3=1\).  The middle diagonal is
the product of the split scalar \(HI\) and \(a(-H^{-2})\), again
without a cocycle.  Thus
\begin{equation}\label{eq:app-bruhat-lift}
 wn(H)=n(H^{-1})a(-H^{-2})\bar n(H^{-1})
\end{equation}
on the quotient.  If \(H^{-1}\in\mathcal O\), the last factor is in
the compact splitting.  The latter agrees with the sections on unit
diagonals, \(w\), and integral upper root elements by
\cite[Proposition~0.1.3, p.~44]{KP84}; conjugation by \(w\) gives
agreement on the lower root group.  The diagonal \(a(-1)\) can
therefore be absorbed in a compact factor in any spherical
calculation.

Let \(\sigma^-\) be the unramified representation of the covering
torus with central character the Weyl conjugate of
\eqref{eq:theta-central-datum}.  Its norm exponent on \(a(y)\) is
\(-1/6\).  Modulo scalars, the subgroup of torus elements of
valuation divisible by three is maximal abelian: the tame group
\(F_{\mathfrak p}^\times/F_{\mathfrak p}^{\times3}\) has one
valuation and one unit coordinate over \(\Z/3\Z\), and the
commutator pairing pairs those coordinates nondegenerately.
The torus construction in \cite[Section~0.3, pp.~54--55, and
Section~I.1, p.~58]{KP84} realizes \(\sigma^-\) by induction from a
character of this subgroup.  Its index is three, so the unramified
Heisenberg representation has dimension three.  The canonical
unramified extension is trivial on unit lifts
\cite[Section~I.1, p.~60]{KP84}; the induced basis vector supported
on that subgroup is therefore unit fixed.  Call it \(e_0\), and set
\begin{equation}\label{eq:app-Heisenberg-basis}
 e_j=Q^{-j/6}\sigma^-(a(\varpi^j))e_0,\qquad j\in\Z.
\end{equation}
The central element \(a(\varpi^3)\) acts by \(Q^{1/2}\), so
\(e_{j+3}=e_j\).  These three vectors form a basis.  From
\eqref{eq:app-torus-laws} and unit invariance one obtains
\begin{equation}\label{eq:app-torus-on-e0}
 \sigma^-(a(\varpi^j u))e_0
   =Q^{j/6}\chi(u)^{2j}e_j
       \qquad(j\in\Z,\ u\in\mathcal O^\times).
\end{equation}

Use the vector-valued realization of normalized induction
\[
 I^-=\operatorname{Ind}_{\widetilde B}^{\widetilde{\PGL}_2}
                   (\sigma^-),\qquad
 f(n(x)a(y)g)=|y|_{\mathfrak p}^{1/2}
                 \sigma^-(a(y))f(g).
\]
Its spherical vector \(f^\circ\) is determined by
\(f^\circ(k)=e_0\) on the compact subgroup.  The total norm exponent
in this transformation law is
\(1/2-1/6=1/3\).  The compact-invariant line of this normalized
unramified induction is one dimensional by
\cite[Lemma~I.1.3, p.~60]{KP84}.  By
\cite[Theorem~I.2.9(a),(f), p.~72]{KP84}, the local exceptional
representation is the irreducible subrepresentation of \(I^-\) and has
a nonzero compact-invariant vector.  That vector is a scalar multiple of
\(f^\circ\), so the subrepresentation contains \(f^\circ\).

\subsection{The shells of the Jacquet integral}

For a smooth \(f\in I^-\), define the vector Jacquet functional by
\[
 \mathscr J(f)=\lim_{L\to\infty}
   \int_{\varpi^{-L}\mathcal O}f(wn(x))\overline{\psi(x)}\,\dd x.
\]
This limit stabilizes.  Indeed, when \(|x|\) is sufficiently large,
the right invariance of \(f\) makes \(\bar n(x^{-1})\) fix \(f\).
The lifted Bruhat formula then expresses \(f(wn(x))\) as
\[
 |x|_{\mathfrak p}^{-1}
       \sigma^-(a(-x^{-2}))f(1).
\]
Choose \(m\) so that \(1+\varpi^m\mathcal O\) consists of cubes and
its unit lifts fix the vector \(f(1)\) in the smooth torus
representation.  On a shell \(x=\varpi^{-l}v\), the displayed
expression is unchanged when \(v\) is replaced by
\(vu\), \(u\in1+\varpi^m\mathcal O\): the norm is unchanged, all
cocycles with \(u\) are trivial, and \(a(u^{-2})\) fixes \(f(1)\).
Each such unit coset is an additive coset of
\(\varpi^m\mathcal O\).  For \(l>m\), the character
\(\psi(\varpi^{-l}v)\) has zero integral on it.  All sufficiently
large shells therefore vanish.  This argument applies to every
smooth \(f\), including each of its right translates.

If \(t\in F_{\mathfrak p}\), choose a ball large enough to contain
\(t\) and to realize the stable values for both \(f\) and
\(\pi(n(t))f\).  Translation \(x\mapsto x+t\) preserves that ball
and gives
\(\mathscr J(\pi(n(t))f)=\psi(t)\mathscr J(f)\).
Thus each scalar coordinate of \(\mathscr J\) is a Whittaker functional.
This is a direct stable version, at the present inducing parameter,
of the Jacquet functionals introduced in
\cite[Section~I.3, pp.~74--75]{KP84}.

For \(e\ge0\) put \(J_e=\mathscr J(\pi(a(\varpi^e))f^\circ)\); explicitly,
\begin{equation}\label{eq:app-Jacquet}
 J_e=\int_{F_{\mathfrak p}}
        f^\circ(wn(x)a(\varpi^e))\overline{\psi(x)}\,\dd x.
\end{equation}
The integral retains the stable interpretation just established.
We now evaluate it shell by shell.

Put \(x=\varpi^eH\) and use the first two lift identities above.
The measure contributes \(Q^{-e}\), the modular factor of
\(a(\varpi^{-e})\) contributes \(Q^{e/2}\), and its torus norm
factor contributes \(Q^{-e/6}\).  Thus every shell has the common
factor
\begin{equation}\label{eq:app-common-factor}
 Q^{-e+e/2-e/6}=Q^{-2e/3}.
\end{equation}
If \(H\in\mathcal O\), then \(wn(H)\) belongs to the compact
subgroup, and \(\psi(\varpi^eH)=1\).  This region contributes
\(e_{-e}\) inside the common factor.

Suppose instead that \(H=\varpi^{-l}v\), with \(l\ge1\) and
\(v\in\mathcal O^\times\).  In \eqref{eq:app-bruhat-lift} the last
factor is compact, while the first acts trivially in the inducing
transformation law.  The modular factor of \(a(-H^{-2})\) is
\(Q^{-l}\), and its torus norm factor is \(Q^{l/3}\).  Furthermore,
\[
 a(\varpi^{2l}v^{-2})
   =\mathcal C(\varpi^{2l},v^{-2})^{-1}
       a(\varpi^{2l})a(v^{-2}),\qquad
 \mathcal C(\varpi^{2l},v^{-2})^{-1}=\chi(v)^{-8l}.
\]
The minus sign in \(-H^{-2}\) has no effect.  Multiplication by
\(a(\varpi^{-e})\) changes the basis index by \(-e\), without a
cocycle between the pure uniformizer powers.  Finally
\(\dd H=Q^l\,\dd v\).  It follows that this shell contributes,
inside \eqref{eq:app-common-factor},
\begin{equation}\label{eq:app-shell}
 Q^{l/3}e_{-e+2l}
       \int_{\mathcal O^\times}
         \chi(v)^{-8l}\overline{\psi(\varpi^{e-l}v)}\,\dd v.
\end{equation}
This accounts separately for the measure, modular, norm, and
Hilbert-symbol factors.

Define the normalized prime Gauss sums
\[
 g_j=Q^{-1/2}\sum_{x\bmod\mathfrak p}
          \chi(x)^j\psi(x/\varpi),\qquad j\in\Z.
\]
The powers in these sums retain the unit mask.  Since
\(\chi(-1)=1\), replacing \(\psi(x/\varpi)\) by its complex
conjugate has no effect on \(g_j\).  For \(j\not\equiv0\pmod6\),
the usual finite-field Gauss identity gives
\[
 |g_j|=1,\qquad g_jg_{-j}=\chi(-1)^j=1.
\]
In particular \(g_2g_4=1\).  These identities follow, for example,
by expanding the product of the two sums, substituting \(x=ty\),
and using the additive sum in \(y\); the result is
\(Q\chi(-1)^j\) before division by \(Q\).

Write \(I_{e,l}\) for the integral in \eqref{eq:app-shell}.  The
character there is \(\chi^{-2l}\), since \(-8l\equiv-2l\pmod6\).
Its exact values are
\begin{equation}\label{eq:app-shell-integrals}
 I_{e,l}=
 \begin{cases}
  1-Q^{-1},&1\le l\le e,\quad 3\mid l,\\
  0,&1\le l\le e,\quad 3\nmid l,\\
  Q^{-1/2}g_{-2l},&l=e+1,\quad 3\nmid l,\\
  -Q^{-1},&l=e+1,\quad 3\mid l,\\
  0,&l>e+1.
 \end{cases}
\end{equation}
Indeed, for \(l\le e\) the additive character is trivial, so unit
orthogonality gives the first two cases.  For \(l=e+1\), integration
over the residue classes gives \(Q^{-1}\) times the corresponding
prime Gauss sum; the trivial character with its unit mask gives
\(\sum_{x\ne0}\psi(x/\varpi)=-1\).  For \(l>e+1\), split each unit
class modulo \(\varpi^{l-e-1}\) into its \(Q\) lifts modulo
\(\varpi^{l-e}\).  The multiplicative character is unchanged on
these lifts, whereas the additive factors sum to
\(\sum_{t\bmod\mathfrak p}\psi(t/\varpi)=0\).
This also proves the asserted stabilization of
\eqref{eq:app-Jacquet}.

We can now sum the shells.  For \(k\ge0\), the integral unit ball
and the shells \(l=3,6,\ldots,3k\) have the same basis index and
combined coefficient
\[
 1+(1-Q^{-1})\sum_{j=1}^{k}Q^j=Q^k.
\]
For \(e=3k\), the terminal shell \(l=3k+1\) has coefficient
\(Q^{k-1/6}g_4\) and basis vector \(e_2\).  For \(e=3k+1\), the
terminal shell \(l=3k+2\) has coefficient
\(Q^{k+1/6}g_2\) and basis vector \(e_0\).  For \(e=3k+2\), the
terminal shell \(l=3k+3\) has coefficient \(-Q^k\) and the same
basis vector \(e_1\) as the preceding combined contribution.
Multiplying by \eqref{eq:app-common-factor} therefore gives
\begin{align}
 J_{3k}
   &=Q^{-k}(e_0+Q^{-1/6}g_4e_2),\notag\\
 J_{3k+1}
   &=Q^{-k}(Q^{-1/2}g_2e_0+Q^{-2/3}e_2)
     =Q^{-k-1/2}g_2(e_0+Q^{-1/6}g_4e_2),\notag\\
 J_{3k+2}&=0.                                      \label{eq:app-J-values}
\end{align}
The second equality uses \(g_2g_4=1\).

In particular \(J_0=e_0+Q^{-1/6}g_4e_2\ne0\).  Choose a scalar
coordinate nonzero on \(J_0\).  Its Jacquet functional is nonzero on
the exceptional subrepresentation, since that subrepresentation
contains \(f^\circ\).  The local Whittaker space is one dimensional
by \cite[Corollary~I.3.6, p.~79]{KP84}.  Hence every spherical Whittaker
function normalized at one has the ratios in
\eqref{eq:app-J-values}:
\begin{equation}\label{eq:app-pure-values}
 W(a(\varpi^{3k}))=Q^{-k},\qquad
 W(a(\varpi^{3k+1}))=Q^{-k-1/2}g_2,\qquad
 W(a(\varpi^{3k+2}))=0.
\end{equation}
The middle term is
\(Q^{-k-1}\sum_x\chi(x)^2\psi(x/\varpi)\), as claimed.
This computation agrees with the central-step statement of
\cite[Theorem~I.4.3, p.~87]{KP84}: on \(a(\varpi^3)\) the Weyl-conjugate
inducing character is \(Q^{1/2}\) and the modular half-character is
\(Q^{-3/2}\), so their product is \(Q^{-1}\).

\subsection{Negative valuations and the unit normalization}

For \(e<0\), choose \(x\in\mathcal O\) for which
\(\psi(\varpi^e x)\ne1\).  The spherical vector is fixed by
\(n(x)\), while the split root relation and Whittaker equivariance
give
\[
 W(a(\varpi^e))
 =W(a(\varpi^e)n(x))
 =\psi(\varpi^e x)W(a(\varpi^e)).
\]
Thus the value is zero.  For a unit \(u\), the compact section of
\(a(u)\) fixes the spherical vector.  The torus law gives
\[
 a(\varpi^e u)=\mathcal C(\varpi^e,u)^{-1}a(\varpi^e)a(u),
\]
where \(\mathcal C\) includes the scalar action on a genuine vector.
Consequently
\begin{equation}\label{eq:app-unit-rule}
 W(a(\varpi^e u))=\chi(u)^{2e}W(a(\varpi^e)).
\end{equation}
Together with \eqref{eq:app-pure-values}, this proves
\eqref{eq:theta-local-values}.

There is a useful check that fixes the orientation without referring
to a convention for the name of a cubic Gauss sum.  At a tame place
the residue character is unchanged when \(\varpi\) is replaced by
\(\varpi u\).  Substitution \(x=uv\) in the residue sum yields
\[
 Q^{-1}\sum_{x\bmod\mathfrak p}\chi(x)^2\psi(x/(\varpi u))
 =\chi(u)^2 Q^{-1}\sum_{v\bmod\mathfrak p}
                         \chi(v)^2\psi(v/\varpi).
\]
This is exactly the \(e=1\) transformation in
\eqref{eq:app-unit-rule}.  Thus the Gauss phase and the unit factor
are a single compatible normalization.  Changing the sign of an
additive Gauss argument is harmless because \(\chi(-1)=1\).
